\documentclass{homework}
% \usepackage{graphicx} % Required for inserting images
\usepackage{amsmath}
\usepackage{gensymb}
\usepackage{subfig}
\usepackage{float}

\class{ECEN/MCEN 5448}
\title{Homework 1}
\author{Sean Jennings}
\date{September 18, 2023}

\begin{document}
\maketitle

\section{State space model and Linearization}
\subsection{(Vehicle Platooning)}
For truck $i\in\{1, 2\}$, the position of the truck is given by $x_i$ and the
input is given by $v_i$. The dynamics of each truck $i$ satisfies:
\begin{equation}
    \ddot{x}_i = v_i
\end{equation}
where $u := v_1$, so that the input to the first truck is defined as the input
to the connected system, and the controller
\begin{equation}
    v_2 = K_p(x_2 - x_1) + K_v (\dot{x}_2 - \dot{x}_1)
\end{equation}
governs the input to the second truck, where $K_p,K_v \in \R$ are controller
constants.

We define the state ${\bf x}$ to be:
\begin{equation}
    {\bf x} := [x_1, \dot{x_1}, x_2, \dot{x_2}]^T
\end{equation}
Then, the derivative is given by:
\begin{equation}
    {\bf \dot{x}} = \begin{bmatrix}
        \dot{x}_1, \\
        \ddot{x}_1, \\
        \dot{x}_2, \\
        \ddot{x}_2
    \end{bmatrix} = \begin{bmatrix}
        \dot{x}_1, \\
        u, \\
        \dot{x}_2, \\
        K_p(x_2 - x_1) + K_v (\dot{x}_1 - \dot{x}_2)
    \end{bmatrix}
\end{equation}
and in matrix form, we have
\begin{equation}
    {\bf \dot{x}} = A {\bf x} + B {\bf u}
\end{equation}
where
\begin{equation}
    A = \begin{bmatrix}
        0 & 1 & 0 & 0 \\
        0 & 0 & 0 & 0 \\
        0 & 0 & 0 & 1 \\
        -K_p & -K_v & K_p & K_v
    \end{bmatrix}, \qquad B = \begin{bmatrix}
        0 \\
        1 \\
        0 \\
        0
    \end{bmatrix}.
\end{equation}
Since our observation is given by 
\begin{equation}
    {\bf y} := [x_1, x_2]^T = \begin{bmatrix}
        1 & 0 & 0 & 0 \\
        0 & 0 & 1 & 0
    \end{bmatrix} {\bf x},
\end{equation}
in matrix form ${\bf y} = C {\bf x} + D {\bf u}$, we have
\begin{equation}
    C = \begin{bmatrix}
        1 & 0 & 0 & 0 \\
        0 & 0 & 1 & 0
    \end{bmatrix}, \qquad D = {\bf 0}_{2 \times 1}.
\end{equation}
Therefore, the platooning system is an LTI system with input ${\bf u}$, state ${\bf x}$,
observation ${\bf y}$ and matrices $A$, $B$, $C$, $D$ defined as above.

\subsection{(Unicycle Model)}

We designate state and input vectors to be
\begin{equation}
    {\bf x} := \begin{bmatrix}
        p_x \\
        p_y \\
        \theta \\
    \end{bmatrix}, \qquad {\bf u} := \begin{bmatrix}
        v \\ \omega
    \end{bmatrix}
\end{equation}
and thus ${\bf \dot{x}} = f({\bf x}, {\bf u})$ where $f$ is
\begin{equation}
    f({\bf x}, {\bf u}) = \begin{bmatrix}
        {\bf u}[0] \cos ({\bf x}[2]) \\
        {\bf u}[0] \sin ({\bf x}[2]) \\
        {\bf u}[1]
    \end{bmatrix}.
\end{equation}
Note that ${\bf x}[i]$ denotes $i$th index of vector ${\bf x}$ and is zero-indexed.
For instance, ${\bf u}[0] = v$.

\begin{letterenum}
    \item To compute the linearization about ${\bf x}^* = {\bf 0}_3$ and
    ${\bf u}^* = {\bf 0}_2$, we define
    \begin{equation}
        \delta x = {\bf x} - {\bf x}^* = {\bf x}, \qquad
        \delta u = {\bf u} - {\bf u}^* = {\bf u}
    \end{equation}
    Since
    \begin{equation}
        \frac{d (\delta x)}{dt} = \frac{d}{dt} ({\bf x} - {\bf x}^*)
        = f({\bf x}, {\bf u}),
    \end{equation}
    a Taylor Series expansion of $f$ about $({\bf x}^*, {\bf u}^*)$ yields
    \begin{equation}
        \frac{d (\delta x)}{dt} = f(\xc^*, \uc^*)
            + \frac{\delta f}{\delta \xc}(\xc^*, \uc^*) \delta x
            + \frac{\delta f}{\delta \uc}(\xc^*, \uc^*) \delta u
            + \mathcal{O}(\delta x^2) + \mathcal{O}(\delta u^2).
    \end{equation}
    Thus, we compute the jacobian of $f$ with respect to $\xc$:
    \begin{equation}
        \frac{\delta f}{\delta \xc}(\xc, \uc) = \begin{bmatrix}
            0 & 0 & - \uc[0] \sin ( \xc[2]) \\
            0 & 0 & \uc[0] \cos ( \xc[2]) \\
            0 & 0 & 0
        \end{bmatrix}
        \label{eq:jacobian-x}
    \end{equation}
    and w.r.t $\uc$:
    \begin{equation}
        \frac{\delta f}{\delta \uc}(\xc, \uc) = \begin{bmatrix}
            \cos ( \xc[2]) & 0 \\
            \sin ( \xc[2]) & 0 \\
            0 & 1
        \end{bmatrix}.
        \label{eq:jacobian-u}
    \end{equation}
    Evaluating at $(\xc^*, \uc^*)$ we have
    \begin{align}
        A := \frac{\delta f}{\delta \xc}(\xc, \uc) = {\bf 0}_{3 \times 3} \\
        B := \frac{\delta f}{\delta \uc}(\xc, \uc) = \begin{bmatrix}
            1 & 0 \\
            0 & 0 \\
            0 & 1
        \end{bmatrix}
    \end{align}
    so that our linearization is given by:
    \begin{equation}
        \frac{d (\delta x)}{dt} \approx B \delta u
    \end{equation}

    \item To show that the trajectory $\omega(t) = v(t) = 1$, $p_x(t) = \sin(t)$,
    $p_y(t) = 1 - \cos(t)$, and $\theta(t) = t$ we take the derivatives of each curve
    \begin{align}
        \dot{p}_x(t) = \frac{d}{dt}(\sin(t)) &= \cos(t) = v(t) \cos(\theta(t)) \\
        \dot{p}_y(t) = \frac{d}{dt}(1 - \cos(t)) &= \sin(t) = v(t) \sin(\theta(t)) \\
        \dot{\theta}(t) = \frac{d}{dt}(t) &= 1 = \omega(t).
    \end{align}
    and see that the system's differential equations are satisfied. Then,
    to find the linearization about the trajectory we define the change of variables:
    \begin{equation}
        \delta x(t) := \xc(t) - \xc_{traj}(t), \qquad \delta u(t) := \uc(t) - \uc_{traj}(t)
    \end{equation}
    Then, the Taylor expansion of the derivative of $\delta x$ is given by:
    \begin{align}
        \frac{d(\delta x)}{dt} &= \frac{d\xc}{dt} - \frac{d\xc_{traj}}{dt} \\
            &= \biggl(f(\xc_{traj}, \uc_{traj})
                + \frac{\delta f}{\delta \xc} (\xc_{traj}, \uc_{traj}) \delta x
                + \frac{\delta f}{\delta \uc} (\xc_{traj}, \uc_{traj}) \delta u
                + \text{H.O.T.}
                \biggr) \\ & \qquad - f(\xc_{traj}, \uc_{traj}) \\
            &\approx \frac{\delta f}{\delta \xc} (\xc_{traj}, \uc_{traj}) \delta x
                + \frac{\delta f}{\delta \uc} (\xc_{traj}, \uc_{traj}) \delta u
    \end{align}
    where the dependence of $\delta x$, $\xc_{traj}$, $\uc_{traj}$ upon time has
    been omitted for brevity.
    Substituting the trajectory curves into the expression for the jacobian w.r.t. 
    $\xc$ (equation (\ref{eq:jacobian-x})) and
    w.r.t. $\uc$ (equation (\ref{eq:jacobian-u})), we obtain the final result
    for the linearization about the trajectory:
    \begin{equation}
        \frac{d(\delta x(t))}{dt} = \begin{bmatrix}
            0 & 0 & -\sin(t) \\
            0 & 0 & \cos(t) \\
            0 & 0 & 0
        \end{bmatrix} \delta x + \begin{bmatrix}
            \cos(t) & 0 \\
            \sin(t) & 0 \\
            0 & 1
        \end{bmatrix} \delta u
    \end{equation}
    We conclude by noting that the linearization about this particular trajectory
    is a LTV system.
\end{letterenum}

\section{Linear Algebra}
\subsection{}
\begin{letterenum}
    \item Let $\alpha, \beta, \gamma$ be elements of the set $\R^{m \times n}$ with
        $m, n \in \N$. Let $a_{ij}$, $b_{ij}$, $c_{ij}$ denote the $i$th $j$th entry of 
        the matrices $\alpha$, $\beta$, $\gamma$, respectively. Let $r,s \in \R$ be
        scalars. 

        The usual matrix addition defines the $ij$th entry of matrix $\alpha + \beta$
        as follows
        \begin{equation}
            (\alpha + \beta)_{ij} = a_{ij} + b_{ij}
        \end{equation}

        The scalar multiplication operation defines the $ij$th entry of matrix
        of $r\alpha$ to be:
        \begin{equation}
            (r\alpha)_{ij} = r a_{ij}
        \end{equation}

        We prove that the set $V = (\R_{m\times n}, +, \cdot)$ over the field of
        $\R$ is a vector space. In the following, we will use the fact that two matrices 
        $\alpha$ and $\beta$ are equal if and only if each of their entries are equal, i.e.
        $a_{ij} = b_{ij}$ for all $1 \leq i \leq m$, $1 \leq j \leq n$. All arguments involving
        indices $i,j$ are assumed to hold for all valid pairs.

        \begin{itemize}
            \item Since $a_{ij} + b_{ij} \in \R$, each element
                $(\alpha + \beta)_{ij} \in \R$ and
                $\alpha + \beta \in \R^{m \times n}$. Thus, $V$ is closed under addition.
            \item Commutativity and associativity of addition follow similarly
                from respective properties of $\R$.
            \item The matrix ${\bf 0}_{m \times n} \in \R_{m \times n}$ which
                has $({\bf 0}_{m \times n})_{ij} = 0$ for all $i,j$ is clearly
                the zero vector since $(\alpha + 0)_{ij} = a_{ij} + 0 = a_{ij}$.
            \item The additive inverse $-\alpha$ has entries $(-\alpha)_{ij} = -a_{ij}$.
                for all $i,j$. Since $-a_{ij} \in \R$ for all $i,j$, we have
                    $-\alpha \in \R^{m \times n}$. Clearly, $\alpha + (-\alpha) = 0_{m \times n}$.
            \item For every $r \in \R$, $(r\alpha)_{ij} = r a_{ij} \in \R$.
                So $r\alpha \in \R^{m \times n}$ and $V$ is closed under scalar
                multiplication.
            \item Clearly, $1 \alpha = \alpha$ since $(1\alpha)_{ij} = 1 a_{ij} = a_{ij}$.
            \item $[(rs)\alpha]_{ij} = (rs)a_{ij} = r(s a_{ij}) = [r (s \alpha)]_{ij}$
                follows from associativity of multiplication on $\R$.
            \item $[r (\alpha + \beta)]_{ij} = r(a_{ij} + b_{ij}) = r a_{ij} + r b_{ij} = [r \alpha]_{ij} + [r \beta]_{ij}$
            \item $[(r+s) \alpha]_{ij} = (r+s)a_{ij} = r a_{ij} + s a_{ij} = [r \alpha]_{ij} + [s \alpha]_{ij}$
        \end{itemize}
        The final two properties follow from the distributive law of the field $\R$.

        Thus, the set $\R^{m \times n}$ with operations $+$ and $\cdot$ over the field $\R$
        satisfies all the properties of a vector space. Clearly, $\R^{2 \times 2}$
        is a special case of $\R^{m \times n}$ and thus is also a vector space.

    \item The dimension of $\R^{m \times n}$ is $m \cdot n$ so the dimension of
        $\R^{2 \times 2}$ is 4.

    \item Define $\epsilon_{IJ}$ as the matrix with the $ij$th entry containing:
        \begin{equation}
            (\epsilon_{IJ})_{ij} = \biggl\{
                \begin{array}{ll}
                    1 &\qquad \text{if }I=i\text{ and }J=j \\
                    0 &\qquad \text{otherwise}
                \end{array}
            \biggr.
        \end{equation}
        Then, the set of vectors $\mathcal{B} = \{\epsilon_{IJ}: I,J \in \N, 1 \leq I \leq m, 1 \leq J \leq n\}$
        form a basis for $\R^{m \times n}$. To see this is a basis, we first verify
        that it spans $\R^{m \times n}$. Let $\alpha = \sum_{i=1}^{m} \sum_{j=1}^{n} c_{ij} \epsilon_{ij}$,
        be a linear combination of the basis $\mathcal{B}$, with scalars $c_{ij} \in \R$.
        By selecting scalars such that $c_{ij} = (\alpha)_{ij}$, we can see the previous equality is
        satified and $\alpha$ is indeed in the span of $\mathcal{B}$.

        To check linear indepedence, examining the $ij$th entry of $\beta$, a linear combination
        of $\mathcal{B}$
        \begin{equation}
            \beta := (\sum_{I=1}^{m}) \sum_{J=1}^{n} c_{IJ} \epsilon_{IJ} = c_{IJ}
        \end{equation}
        Thus, $\beta = 0$ if and only if $c_{IJ} = 0$ for all $I,J$,
        $1 \leq I \leq m$, $1 \leq J \leq n$. This implies linear independence and
        we have now proven $\mathcal{B}$ is a basis for $\R^{m \times n}$.

        In particular, for $\R^{2 \times 2}$ our suggested basis is:
        \begin{equation}
            \epsilon_{11} = \begin{bmatrix}
                1 & 0 \\ 0 & 0
            \end{bmatrix} \qquad
            \epsilon_{12} = \begin{bmatrix}
                0 & 1 \\ 0 & 0
            \end{bmatrix} \qquad
            \epsilon_{21} = \begin{bmatrix}
                0 & 0 \\ 1 & 0
            \end{bmatrix} \qquad
            \epsilon_{22} = \begin{bmatrix}
                0 & 0 \\ 0 & 1
            \end{bmatrix} \qquad
        \end{equation}


\end{letterenum}

\subsection{}
The Cayley-Hamilton Theorem says that for any $n \times n$ square matrix, $B$,
the constants $b_1, b_2, \dots, b_n$ of the characteric polynomial satisfy
the equation:
\begin{equation}
    \Delta(B) = B^n + b_1 B^{n-1} + b_2 B^{n-2} + \dots + b_{n-1} B + b_n I_{n\times n} = 0_{n\times n}
\end{equation}
Since the characteristic polynomial, $\Delta$, for our particular $A$ is
\begin{equation}
    \Delta(\lambda) = (1 - \lambda) (2 - \lambda) - 0 = \lambda^2 - 3 \lambda + 2;
\end{equation}
by Cayley-Hamilton, we have $A^2 + (-3) A + 2 I_{n \times n}=0$. We note
the left-hand side is a linear combination of $\{I, A, A^2\}$, equal to
0, and has non-trivial coefficients. Thus, this set of vectors is linearly dependent. 

Alternatively, without using the Cayley-Hamilton, we can verify linear dependence by hand:

\begin{equation}
    A^2 = \begin{bmatrix}
        1 & 1 \\ 0 & 2
    \end{bmatrix} \begin{bmatrix}
        1 & 1 \\ 0 & 2
    \end{bmatrix} = \begin{bmatrix}
        1 & 3 \\ 0 & 4
    \end{bmatrix}
\end{equation}
To check whether the set $\{I, A , A^2\}$ is linearly dependent we look for
non-trivial solutions to the equation $aI + bA + cA^2 = 0$ where $a,b,c\in \R$,
which yields the system of equations:
\begin{align}
    a + b + c &= 0 \\
        b + 3c &= 0 \\
    a + 2b + 4c &= 0.
\end{align}
We then convert to a matrix representation
\begin{equation}
    \begin{bmatrix}
        1 & 1 & 1 \\
        0 & 1 & 3 \\
        1 & 2 & 4 \\
    \end{bmatrix} \begin{bmatrix}
        a \\ b \\ c
    \end{bmatrix} = 0 \\
\end{equation}
and row-reduce:
\begin{equation}
    \begin{bmatrix}
        1 & 0 & -2 \\
        0 & 1 & 3 \\
        0 & 0 & 0
    \end{bmatrix} \begin{bmatrix}
        a \\ b \\ c
    \end{bmatrix} = 0
\end{equation}

So for any $c \in \R$, contants $a = 2c$, $c = -3b$ gives a non-trivial solution.
to the equation $aI + bA + cA^2 = 0$. For instance, $a=2$, $b=-3$, $c=1$, gives
the linear combination:
\begin{equation}
    2 \begin{bmatrix}
        1 & 0 \\ 0 & 1
    \end{bmatrix} + (-3) \begin{bmatrix}
        1 & 1 \\ 0 & 2
    \end{bmatrix} + \begin{bmatrix}
        1 & 3 \\ 0 & 4
    \end{bmatrix} = \begin{bmatrix}
        2 - 3 + 1 & -3 + 3 \\
        0 & 2 - 6 + 4
    \end{bmatrix} = 0.
\end{equation}

Therefore, the set $\{I, A, A^2\}$ is linearly dependent.
\newcommand{\Nc}{\mathcal{N}}
\newcommand{\basis}{\mathcal{B}}

\subsection{}
\begin{letterenum}
    \item We rearrange the given equation for $(T(\theta))_i$ so that
        the matrix form is particularly apparent:
        \begin{align}
            (T(\theta))_i = u_i &= \frac{1}{|\Nc_i} \sum_{j \in \Nc_i} \theta_j - \theta_i \\
                &= \frac{1}{|\Nc_i|} \bigr[ \sum_{j \in \Nc_i} \theta_j - \sum_{j \in \Nc_i} \theta_i] \\
                &= \frac{1}{|\Nc_i| \sum_{j \in \Nc_i}} - \theta_i \\
                &= \sum_{j=1}^m a_{ij}\theta_i
                \label{matrix form}
        \end{align}
        where 
        \begin{equation}
            a_{ij} = \biggl\{
                \begin{array}{ll}
                    -1 &\qquad \text{if }i=j, \\
                    1/|\Nc_i| & \qquad \text{if }j |\in \Nc_i \\
                    0 &\qquad \text{otherwise}.
                \end{array}
            \biggr.
        \end{equation}
        Clearly, $T$ can be represented as a matrix operation $T(\theta) = A \theta$,
        where $A$'s $ij$th entry is given by $a_{ij}$, and is therefore a linear
        transformation.

        For our particular set of neighbors, our matrix representation of the transform $A$
        is
        \begin{equation}
            A = \begin{bmatrix}
                -1 & 1 & 0 & 0 \\
                1/2 & -1 & 1/2 & 0 \\
                0 & 1/2 & -1 & 1/2 \\
                0 & 0 & 1 & -1
            \end{bmatrix}.
            \label{eq:matrix-form-written}
        \end{equation}
        Reducing to row echelon form
        \begin{equation}
            \begin{bmatrix}
                1 & -1 & 0 & 0 \\
                0 & 1 & -1 & 0 \\
                0 & 0 & 1 & -1 \\
                0 & 0 & 0 & 0
            \end{bmatrix}
        \end{equation}
        shows that $\mathcal{B}_N = \{[1, 1, 1, 1]^T\}$
        is a basis for the null space.

        Since $\epsilon_1$, $\epsilon_2$, and $\epsilon_3$ (where $\epsilon_i$ comes from
        the canonical basis for $\R^4$) extend $\mathcal{B}_N$ to form a basis
        for $\R^4$, the set $\mathcal{B}_R = \{T\epsilon_1, T\epsilon_2, T\epsilon_3\}$ forms
        a basis for the range space. That is,
        \begin{equation}
            \mathcal{B}_R = \begin{Bmatrix}
                \begin{bmatrix}
                    -1 \\ 1/2 \\ 0 \\ 0
                \end{bmatrix}, \begin{bmatrix}
                    1 \\ -1 \\ 1/2 \\ 0
                \end{bmatrix}, \begin{bmatrix}
                    0 \\ 1/2 \\ -1 \\ 1
                \end{bmatrix}
            \end{Bmatrix}
        \end{equation}.
        \item The matrix representation in the standard basis is given by
        equation (\ref{eq:matrix-form-written}).


        \item 
        Let $\basis_1 = \{\epsilon_1, \epsilon_2, \epsilon_3, \epsilon_4\}$ be the standard basis
        and $\basis_2$ be the basis given by
        $\{\epsilon_1 - \epsilon_2, \epsilon_2 - \epsilon_3,
            \epsilon_3 - \epsilon_4, \epsilon_4\}$.
            We will write $[P]_{\basis_2}^{\basis_1}$ to denote the expression
            of $\basis_1$ with respect to $\basis_2$. In other words, the $j$th column
            of $[P]_{\basis_2}^{\basis_1}$ is the $j$th basis vector of $\basis{1}$
            expressed in terms of $\basis_2$.
            
            Let $V$ be a finite dimensional vector space with dimension
            $n$ and $\alpha \in V$. Then $[\alpha]_{\basis_1} \in \R^n$ shall
            denote the expression of $\alpha$ as coordinates in $\basis_1$.
            We transform between coordinates with via the relation:
            \begin{equation}
                [\alpha]_{\basis_2} = [P]_{\basis_2}^{\basis_1} [\alpha]_{\basis_1}
            \end{equation}

            Let $T: V \to W$ be a linear operator between vector spaces $V$ and $W$,
            and $\basis_1$, $\basis_2$ be bases of $V$ and $\mathcal{C}_1$, $\mathcal{C}_2$
            be bases of $W$.
            We then use the notation $[T]_{\mathcal{C}_1}^{\basis_1}$ to denote
            the matrix representation of the transformation $T$ from a
            coordinate vector expressed in $\basis_1$ to a coordinate vector
            expressed in $\mathcal{C}_1$.

            This notation allows the conveniently memorable notation of the
            linear transform change of basis:
            \begin{equation}
                [T]_{\mathcal{C}_2}^{\basis_2} = [Q]_{\mathcal{C}_2}^{\mathcal{C}_1}
                    [T]_{\mathcal{C}_1}^{\basis_1}
                    [P]_{\basis_1}^{\basis_2}.
            \end{equation}

            We now return to the problem at hand of expressing the operator
            $T$ of our flocking problem in the terms of the basis $\basis_2$.
            The change of basis matrix is
            \begin{equation}
                [P]_{\basis_1}^{\basis_2} = \begin{bmatrix}
                    1 & 0 & 0 & 0 \\
                    -1 & 1 & 0 & 0 \\
                    0 & -1 & 1 & 0 \\
                    0 & 0 & -1 & 1 \\
                \end{bmatrix}.
                \label{eq:basis-1-to-basis-2}
            \end{equation}
            Note that the first column of this matrix represents the first
            basis vector, called $\beta_1$, of $\basis_2$ with respect to $\basis_1$:
            \begin{align}
                \beta_1 &= \sum_{i=1}^{4} a_i \epsilon_i = \epsilon_1 - \epsilon_2 \\
                [\beta_1]_{\basis_1} &= \begin{bmatrix}
                    1 \\ -1 \\ 0 \\ 0
                \end{bmatrix}.
            \end{align}
            The inverse of this matrix is the expression of $\basis_1$ w.r.t.
            $\basis_2$
            \begin{equation}
                [P]_{\basis_2}^{\basis_1} = \begin{bmatrix}
                    1 & 0 & 0 & 0 \\
                    1 & 1 & 0 & 0 \\
                    1 & 1 & 1 & 0 \\
                    1 & 1 & 1 & 1 \\
                \end{bmatrix}
                \label{eq:basis-2-to-basis-1}
            \end{equation}
            Since there is ambiguity about whether the matrix representation
            of $A'$ refers to the matrix representation using the basis
            $\basis_2$ for just the input, just the output, or both, we
            assume that $A'$ refers to the matrix representation of $T$ with
            both input and output expressed in $\basis_2$. In other words,
            $[T]_{\basis_2}^{\basis_2}$.
            \begin{equation}
                [T]_{\basis_2}^{\basis_2} = [P]_{\basis_2}^{\basis_1}
                    [T]_{\basis_1}^{\basis_1}
                    [P]_{\basis_1}^{\basis_2}
            \end{equation}
            Substituting values from equations (\ref{eq:matrix-form-written}),
            (\ref{eq:basis-1-to-basis-2}), and (\ref{eq:basis-2-to-basis-1})
            and performing the matrix multiplication yields:
            \begin{equation}
                [T]_{\basis_2}^{\basis_2} = \begin{bmatrix}
                    -2 &  1 &  0 &  0  \\ 
                    -0.5 & -0.5 &  0.5 &  0  \\
                    -1 &  1 & -1 &  0.5  \\
                    -1 &  0 &  1 & -0.5  \\
                \end{bmatrix}
            \end{equation}

            We will conclude by noting that basis $\basis_2$ differs from
            the one that the author would expect to provide a more intuitive
            representation of the transformation. By selecting the basis
            $\basis_3 = \{\epsilon_1,
                \epsilon_1 + \epsilon_2,
                \epsilon_1 + \epsilon_2 + \epsilon_3,
                \epsilon_1 + \epsilon_2 + \epsilon_3 + \epsilon_4\}$, we find
            change of basis matrices:
            \begin{align}
                [P]_{\basis_1}^{\basis_3} &= \begin{bmatrix}
                    1 & 1 & 1 & 1 \\
                    0 & 1 & 1 & 1 \\
                    0 & 0 & 1 & 1 \\
                    0 & 0 & 0 & 1
                \end{bmatrix} \\
                [P]_{\basis_3}^{\basis_1} &= \begin{bmatrix}
                    1 & -1 & 0 & 0 \\
                    0 & 1 & -1 & 0 \\
                    0 & 0 & 1 & -1 \\
                    0 & 0 & 0 & 1
                \end{bmatrix}
            \end{align}
            Let $\theta \in \R^4$ be the angle of each bird and
            $\delta\theta = [\theta]_{\basis_3} = [P]_{\basis_3}^{\basis_1} \theta$
            be the representation of vector $\theta$ w.r.t. $\basis_3$.
            % \begin{equation}
            %     \delta\theta = [P]_{\basis_3}^{\basis_1} \theta
            % \end{equation}
            Then, the $i$th entry of $\delta\theta$,
            \begin{equation}
                \delta\theta_i = \biggl\{\begin{array}{ll}
                        \theta_i &\qquad \text{if } i = 4 \\
                        \theta_i - \theta_{i + 1} &\qquad \text{otherwise,}
                \end{array}
                \biggr.
            \end{equation}
            represents the difference between adjacent bird's headings for all but
            the last. Given the particular collection of neighbors,
            $\{\mathcal{N}_i: i=1,\dots,4\}$ contains strictly adjacent birds,
            it is easy to imagine how this representation might promote
            intuitive understanding. Indeed,
            \begin{align}
                [T]_{\basis_3}^{\basis_3} &= [P]_{\basis_3}^{\basis_1}
                        [T]_{\basis_1}^{\basis_1}
                        [P]_{\basis_1}^{\basis_3} \\
                    &= \begin{bmatrix}
                        -1.5 & 0.5 & 0 & 0 \\
                        0.5 & -1 & 0.5 & 0 \\
                        0 & 0.5 & -1.5 & 0 \\
                        0 & 0 & 1 & 0 \\
                    \end{bmatrix}
            \end{align}
            has all zeros in the 4th column, indicating that the absolute
            reference point of the flock does not impact the control input,
            only the relative positions.
\end{letterenum}

\section{Matrix Exponentials}

\subsection{}
\begin{letterenum}
    \item 
    \begin{equation}
        A = \begin{bmatrix}
            0 & 1 \\ 0 & 1
        \end{bmatrix}
    \end{equation}
    Characteristic polynomial: $-\lambda(1 - \lambda) = 0$.

    Eigenvalues: $\lambda_1 = 1$, $\lambda_2 = 0$

    Corresponding eigenvectors: $v_1 = [1, 1]^T$, $v_0 = [1, 0]^T$

    $A = PDP^{-1}$ where
    \begin{align}
        P = \begin{bmatrix}
            1 & 1 \\ 1 & 0
        \end{bmatrix} \\
        P^{-1} = \begin{bmatrix}
            0 & 1 \\ 1 & -1
        \end{bmatrix} \\
        D = \begin{bmatrix}
            1 & 0 \\ 0 & 0
        \end{bmatrix} \\
    \end{align}

    Then, the exponential is:
    \begin{align}
        e^{At} &= P e^{Dt} P^{-1} \\
            &= \begin{bmatrix}
                1 & 1 \\ 1 & 0
            \end{bmatrix} \begin{bmatrix}
                e^t & 0 \\ 0 & e^{0t}
            \end{bmatrix} \begin{bmatrix}
                0 & 1 \\ 1 & -1
            \end{bmatrix} \\
            &= \begin{bmatrix}
                1 & e^t - 1 \\ 0 & e^{t}
            \end{bmatrix}
    \end{align}

    \item 
    \begin{equation}
        A = \begin{bmatrix}
            0 & 1 \\ 1 & 0
        \end{bmatrix}
    \end{equation}
    Characteristic polynomial: $\lambda^2 - 1 = 0$.

    Eigenvalues: $\lambda_{1,2} = \pm 1$

    Corresponding eigenvectors: $v_1 = [1, 1]^T$, $v_0 = [1, -1]^T$

    $A = PDP^{-1}$ where
    \begin{align}
        P = \begin{bmatrix}
            1 & 1 \\ 1 & -1
        \end{bmatrix} \\
        P^{-1} = \begin{bmatrix}
            1/2 & 1/2 \\ 1/2 & -1/2
        \end{bmatrix} \\
        D = \begin{bmatrix}
            1 & 0 \\ 0 & -1
        \end{bmatrix} \\
    \end{align}

    Then, the exponential is:
    \begin{align}
        e^{At} &= P e^{Dt} P^{-1} \\
            &= \begin{bmatrix}
                1 & 1 \\ 1 & -1
            \end{bmatrix} \begin{bmatrix}
                e^t & 0 \\ 0 & e^{-t}
            \end{bmatrix} \frac{1}{2} \begin{bmatrix}
                1 & 1 \\ 1 & -1
            \end{bmatrix} \\
            &= \frac{1}{2} \begin{bmatrix}
                e^t + e^{-t} & e^t - e^{-t} \\ e^t - e^{-t} & e^t + e^{-t}
            \end{bmatrix}
    \end{align}

    \item 
    \begin{equation}
        A = \begin{bmatrix}
            0 & 1 \\ -1 & 0
        \end{bmatrix}
    \end{equation}
    Characteristic polynomial: $\lambda^2 + 1 = 0$.

    Eigenvalues: $\lambda_{1,2} = \pm i$

    Corresponding eigenvectors: $v_1 = [1, i]^T$, $v_0 = [i, 1]^T$

    $A = PDP^{-1}$ where
    \begin{align}
        P = \begin{bmatrix}
            1 & i \\ i & 1
        \end{bmatrix} \\
        P^{-1} = \frac{1}{2} \begin{bmatrix}
            1 & -i \\ -i & 1
        \end{bmatrix} \\
        D = \begin{bmatrix}
            i & 0 \\ 0 & -i
        \end{bmatrix} \\
    \end{align}

    Then, the exponential is:
    \begin{align}
        e^{At} &= P e^{Dt} P^{-1} \\
            &= \begin{bmatrix}
                1 & i \\ i & 1
            \end{bmatrix} \begin{bmatrix}
                e^{it} & 0 \\ 0 & e^{-it}
            \end{bmatrix} \frac{1}{2} \begin{bmatrix}
                1 & -i \\ -i & 1
            \end{bmatrix} \\
            &= \frac{1}{2} \begin{bmatrix}
                e^t + e^{-t} & e^t - e^{-t} \\ e^t - e^{-t} & e^t + e^{-t}
            \end{bmatrix}
    \end{align}
    Since $e^{-it} = \cos(t) - i \sin(t)$,
    \begin{align}
        \frac{1}{2} (e^{it} + e{-it}) = \cos(t) \\
        \frac{i}{2} (e^{it} - e{-it}) = -\sin(t) \\
    \end{align}
    and we have a final formula for the exponential:
    \begin{equation}
        e^{At} = \begin{bmatrix}
            \cos(t) & \sin(t) \\
            -\sin(t) & \cos(t)
        \end{bmatrix}
    \end{equation}

\end{letterenum}

\subsection{}
We begin with a theorem proved in lecture: If $A,B \in \R^{n \times n}$ with
$AB = BA$, then
\begin{equation}
    e^{A+B} = e^{B+A} = e^A e^B.
    \label{eq:exp-multiplication-law}
\end{equation}
We note that $A(-A) = (-A)A = -A^2$. Then we have:
\begin{equation}
    e^A e^{-A} = e^{A-A} = e^{0_{n \times n}} =
        \sum_{i=0}^{\infty} \frac{0_{n \times n}^i}{i!} = I_{n \times n}.
\end{equation}

Thus, $(e^A)^{-1} = e^{-A}$.

\subsection{}

Let $A\in\R^{n \times n}$, $n\in\N$.

First, we will prove:
\begin{equation}
    (A^i)^T = (A^T)^i
    \label{eq:matrix-transpose}
\end{equation}
for $i=0,1,2,\dots$.

We proceed by induction. Clearly, $(A^0)^T = I = (A^T)^0$. Now, assume
$(A^i)^T = (A^T)^i$. We want to show, $(A^{i+1})^T = (A^T)^{i+1}$. Using
$(AB)^T = B^T A^T$ for $A,B\in\R^{n\times n}$, we have
\begin{equation}
    (A^{i+1})^T = (A^i A)^T = A^T(A^i)^T = A^T(A^T)^i = (A^T)^{i+1}
\end{equation}
which completes the proof for equation (\ref{eq:matrix-transpose}).

We now use (\ref{eq:matrix-transpose}) to show that
\begin{equation}
    e^{A^T} = (e^A)^T
    \label{eq:exp-transpose}
\end{equation}
Proof:
\begin{equation}
    (e^A)^T = \biggl(\sum_{i=0}^{\infty} \frac{A^i}{i!}\biggr)^T
        = \sum_{i=0}^{\infty} \frac{(A^i)^T}{i!}
        = \sum_{i=0}^{\infty} \frac{(A^T)^i}{i!}
        = e^{A^T}.
\end{equation}

We now complete the proof that if $A$ is skew-symmetric, it's exponential
map is orthogonal.
\begin{equation}
    (e^A)^T e^A = e^{A^T}e^A = e^{A^T + A} = e^{0_{n \times n}} = I_{n\times n}
\end{equation}
where the first equality follows from eq (\ref{eq:exp-transpose}) and the third
equality follows from $A$'s skew-symmetric property. To see that eq
(\ref{eq:exp-multiplication-law}) applies, we verify that $A$ and $A^T$
commute:
\begin{equation}
    AA^T = A(-A) = (-A)A = A^TA
\end{equation}

\section{Harmonic Oscillator}
The non-homogeneous LTI system
\begin{equation}
    \dot{x} = Ax + Bu
\end{equation}
has solution
\begin{align}
    x(t) &= e^{At}x(0) + \int_{0}^{t}e^{A(t-\tau)}Bu(\tau)d\tau \\
        &= e^{At}x(0) + e^{At}
            \biggl( \int_{0}^{t} e^{-A\tau}u(\tau)d\tau \biggr)
            B
\end{align}

\newcommand{\conv}{\int_{0}^{t} e^{-A\tau} d\tau }
\newcommand{\convu}{\int_{0}^{t} e^{-A\tau} u(\tau) d\tau }

Assuming $A$ is invertible, we note:
\begin{equation}
    \conv = -A^{-1}e^{-A\tau} \eval_0^\tau = A^{-1} (I - e^{-A\tau})
\end{equation}

We will use the following shorthand:
\begin{align*}
    c := \cos(t) &\qquad s := \sin(t) \\
    c_{2t} := \cos(2t) &\qquad s_{2t} := \sin(2t)
\end{align*}

Since
\begin{equation}
    A = \begin{bmatrix}
        0 & 1 \\ -1 & 0
    \end{bmatrix},
\end{equation}
we have
\begin{equation}
    A^{-1} = \begin{bmatrix}
        0 & -1 \\ 1 & 1
    \end{bmatrix} \qquad e^{At} = \begin{bmatrix}
        c & s \\ -s & c
    \end{bmatrix}.
\end{equation}


\begin{letterenum}
    \item For $u(t) = 1$, we have:
    \begin{equation}
        \convu = A^{-1}(I - e^{-A\tau})
    \end{equation}
    So that
    \begin{align}
        x(t) &= e^{At}x(0) + A^{-1}(e^{A\tau} - I)B \\
            &= \mat{c & s \\ -s & c} \mat{1 \\ 0} + \mat{0 & -1 \\ 1 & 0}
                \biggl(\mat{c & s \\ -s & c} - I\biggr) \mat{1 \\ 0} \\
            &= \mat{c + s \\ c - s - 1}
    \end{align}
    and
    \begin{equation}
        x(1) = \mat{\cos(1) + \sin(1) \\ \cos(1) - \sin(1) - 1} \approx \mat{1.382 \\ 0.382}
    \end{equation}

    \item For $u(t) = t$, we begin by computing the integral via
    integration by parts:
    \begin{align}
        \int_a^b udv = uv \eval_a^b - \int_{a}^{b} vdu \\
        u = \tau & \qquad v = -A^{-1} e^{-A\tau} \\
        du = d\tau &\qquad dv = e^{-A\tau}d\tau
    \end{align}
    \begin{align}
        \convu &= \int_{0}^{t} \tau e^{-A\tau}d\tau \\
            &= \tau(-A^{-A\tau})\eval_0^t - \int_{0}^{t} (-A^{-1} e^{-A\tau}) d\tau \\
            &= tA^{-1} (I - e^{-At}) + (A^{-1})^2 (I - e^{-At}) \\
            &= A^{-1} (tI + A^{-1}) (I - e^{-At})
    \end{align}
    % This causes Latex to break????
    % so that
    % \begin{align}
    %     x(t) &= e^{At}x(0) + A^{-1}(tI + A^{-1})(e^{At} - I)B \\
    %         &= \mat{c & s \\ -s & c}\mat{1 \\ 0} +
    %             \mat{0 & -1}{1 & 0}\mat{t & -1 \\ 1 & t}\mat{c - 1 \\ -s} \\
    % \end{align}
    % and evaluation at $t=1$ yields:
    % \begin{equation}
        
    % \end{equation}

\end{letterenum}





        
        




\end{document}
